\vfill%
\pagebreak

\section*{Exercises}
\markright{\MakeUppercase{Exercises}}%
\subsection*{Exercise 1}

Write a function \token{reverse(values: [i32])-> [i32]} that returns a new slice
holding the elements of \token{values} in the reverse order. Check it on a slice
of five numbers, on a slice of one, and on the empty slice.

\subsection*{Exercise 2}

What does the following program print? Work it out by hand, then check your
answer by compiling and running it.

\begin{lstlisting}[style=coloredverbatim]
use std::io;

fn main() {
  let s = copy [10, 20, 30, 40, 50];
  let t = (s[1], s[$ - 1], s.len);
  println(t);
  let (a, _, n) = t;
  println(a * cast!i32(n));
  let part = s[1 .. 3] ~ [60];
  for i, x in part {
    print(i, ":", x, " ");
  }
  println();
  match (part.len, part[$ - 1]) {
    (3, 60) => println("three, ending with 60");
    (_, 60) => println("ending with 60");
    _ => println("something else");
  }
}
\end{lstlisting}

\subsection*{Exercise 3}

The following program should replace the first of three readings, which was
wrong, with 12, and raise an alert when the highest reading is above 10. It
contains three errors. Find them, correct them, and compile the program to test
your corrections.

\begin{lstlisting}[style=coloredverbatimError, caption=Ymir program with errors]
use std::io;

fn highest(values: [i32])-> i32? {
  if values.len == 0 {
    return none;
  }
  let mut top = values[0];
  for v in values if v > top {
    top = v;
  }
  top?
}

fn main() {
  let readings: [i32] = [3, 9, 4];
  readings[0] = 12;
  let top = highest(readings);
  if top > 10 {
    println("Alert: a reading of ", top, ".");
  }
}
\end{lstlisting}

\subsection*{Exercise 4}

The function \token{average} of Listing~\ref{lst:collections:average} divides by
0 when the slice is empty. Change it so that it returns an \token{i32?}, with no
value for an empty slice. Then write a function \token{minMax} that returns the
smallest and the largest element of a slice together, or nothing when the slice
is empty. Use both in a program that reads numbers until the user types 0, as
Listing~\ref{lst:collections:numbers} does, and prints their average, their
smallest and their largest. (\textit{Hint: a single \token{match} can take apart
a pair of options.})

\subsection*{Exercise 5}

A palindrome reads the same from left to right and from right to left, such as
\textit{kayak}. Write a function \token{isPalindrome(text: [c8])-> bool}, and a
program that reads a line and says whether it is a palindrome. Compare the first
character with the last, the second with the one before last, and so on. Spaces
count as characters, so \textit{never odd or even} is not a palindrome for this
program.

\subsection*{Exercise 6}

The calendar of Listing~\ref{lst:functions:calendar} used a \token{match} to
name the months and another to count their days. Rewrite \token{monthName} and
\token{daysIn(month, year)} with arrays, as in
Listing~\ref{lst:collections:months}, without forgetting the leap years, and
write a \token{main} that prints the number of days of each month of 2024.

\subsection*{Exercise 7}

Write a program that reads numbers until the user types 0, then prints each
number that was typed more than once, with the number of times it was typed.
(\textit{Hint: use a map.}) In what order are the numbers printed?

\subsection*{Exercise 8}

Pascal's triangle starts with a row holding a single 1. Each next row starts and
ends with 1, and each of its other numbers is the sum of the two numbers above
it: after \token{[1, 3, 3, 1]} comes \token{[1, 4, 6, 4, 1]}. Write a function
\token{nextRow(row: [i32])-> [i32]} that computes the row after \token{row}, and
use it to print the first ten rows of the triangle.

\subsection*{Exercise 9}

A tic-tac-toe board is a grid of 3 by 3 cells, each holding \token{'X'},
\token{'O'}, or \token{'.'} when it is empty. A player wins with three of their
letters on a row, on a column or on one of the two diagonals. Write a function
\token{winner(board: [[c8 ; 3] ; 3])-> c8?} that returns the letter of the
winner, if there is one, and test it on a few boards. (\textit{Hint: a function
that tells whether three cells hold the same letter, and which, is useful eight
times.})

\subsection*{Exercise 10}

Rewrite the function \token{winner} of Exercise~9 for a board stored in a single
array of 9 cells, \token{[c8 ; 9]}, row after row, as in
Section~\ref{sec:collections:flat_grid}. At what index is the cell of row
\token{i} and column \token{j}? Which indices form the first column, and which
the two diagonals? Test the function on the boards of Exercise~9.

\vfill%
\pagebreak

\forceevenpage{}

\section*{Solutions}
\markright{\MakeUppercase{Solutions}}%
\subsection*{Exercise 1}

The loop goes through the indices from the last to the first, and appends each
element to the result. On an empty slice, \token{0 .. 0} is empty, and the
result stays empty.

\lstinputlisting[style=coloredverbatim, caption={Solution for exercise 1, \textit{examples/collections/solutions/reverse.yr}}]{examples/collections/solutions/reverse.yr}
\vspace{-10pt}%
\begin{lstlisting}[style=bashVerb]
[5, 4, 3, 2, 1]
[42]
[]
\end{lstlisting}

\subsection*{Exercise 2}

\begin{lstlisting}[style=bashVerb]
(20, 50, 5)
100
0:20 1:30 2:60
three, ending with 60
\end{lstlisting}

\token{s[\$ - 1]} is the last element, 50, and \token{s.len} is 5, which
\token{cast!i32} turns into an \token{i32} before the product. \token{s[1 .. 3]}
holds the elements at indices 1 and 2, index 3 excluded, and 60 is appended to
it. The indices printed are those of \token{part}, from 0. The pair matched is
\token{(3, 60)}, which the first arm fits; the second would fit it too, but
comes too late.

\subsection*{Exercise 3}

\token{[3, 9, 4]} is an array, and a slice needs a \token{copy}. The first
reading is changed in place, so the slice must be \token{dmut}. And
\token{highest} returns an option, which cannot be compared with 10 before its
value is taken out.

\begin{lstlisting}[style=coloredverbatim, escapechar=@, caption=Solution for exercise 3]
use std::io;

fn highest(values: [i32])-> i32? {
  if values.len == 0 {
    return none;
  }
  let mut top = values[0];
  for v in values if v > top {
    top = v;
  }
  top?
}

fn main() {
  let @\hcb{dmut readings = copy [3, 9, 4]}@;
  readings[0] = 12;
  @\hcb{if let Ok(top) = highest(readings) \{}@
    if top > 10 {
      println("Alert: a reading of ", top, ".");
    }
  @\hcb{\}}@
}
\end{lstlisting}

\subsection*{Exercise 4}

\token{average} and \token{minMax} test the empty slice first, and return
\token{none}. The \token{match} of \token{main} takes apart a tuple of two
options, and the pattern \token{(Ok(avg), Ok((low, high)))} takes apart the
pair inside the second option as well. Both options are empty at the same time,
so the other arm does not have to tell which one is.

\lstinputlisting[style=coloredverbatim, caption={Solution for exercise 4, \textit{examples/collections/solutions/stats.yr}}]{examples/collections/solutions/stats.yr}
\vspace{-10pt}%
\begin{lstlisting}[style=bashVerb, escapechar=@+]
@+\textcolor{teal!80}{alice@dev:\mtilde}@+$ ./stats
Number (0 to stop): 4
Number (0 to stop): 8
Number (0 to stop): 15
Number (0 to stop): 0
Average 9, from 4 to 15.
@+\textcolor{teal!80}{alice@dev:\mtilde}@+$ ./stats
Number (0 to stop): 0
No numbers, no statistics.
\end{lstlisting}

\subsection*{Exercise 5}

The loop needs to go only halfway: the character at index \token{i} is compared
with the one at the same distance from the end, \token{text.len - 1 - i}, and
the second half would compare the same pairs again. The function returns
\token{false} at the first difference.

\lstinputlisting[style=coloredverbatim, caption={Solution for exercise 5, \textit{examples/collections/solutions/palindrome.yr}}]{examples/collections/solutions/palindrome.yr}
\vspace{-10pt}%
\begin{lstlisting}[style=bashVerb, escapechar=@+]
@+\textcolor{teal!80}{alice@dev:\mtilde}@+$ ./palindrome
Text: kayak
"kayak" reads the same both ways.
\end{lstlisting}

\subsection*{Exercise 6}

February is the one month whose length depends on the year, so \token{daysIn}
handles it before reading the array.

\lstinputlisting[style=coloredverbatim, caption={Solution for exercise 6, \textit{examples/collections/solutions/months.yr}}]{examples/collections/solutions/months.yr}
\vspace{-10pt}%
\begin{lstlisting}[style=bashVerb, escapechar=@+]
January 2024: 31 days
February 2024: 29 days
March 2024: 31 days
@+\textit{(9 more lines)}@+
\end{lstlisting}

\subsection*{Exercise 7}

A map from numbers to counts is filled as in
Listing~\ref{lst:collections:letters}. The numbers are printed in the order the
map keeps them, which is neither the order they were typed in nor their
numerical order.

\lstinputlisting[style=coloredverbatim, caption={Solution for exercise 7, \textit{examples/collections/solutions/repeated.yr}}]{examples/collections/solutions/repeated.yr}
\vspace{-10pt}%
\begin{lstlisting}[style=bashVerb, escapechar=@+]
@+\textcolor{teal!80}{alice@dev:\mtilde}@+$ ./repeated
Number (0 to stop): 3
Number (0 to stop): 5
Number (0 to stop): 3
Number (0 to stop): 7
Number (0 to stop): 5
Number (0 to stop): 3
Number (0 to stop): 0
3 was typed 3 times.
5 was typed 2 times.
\end{lstlisting}

\subsection*{Exercise 8}

\token{nextRow} starts the new row with 1, appends the sum of each pair of
neighbours of the row above, and ends with 1. A row of \token{n} numbers has
\token{n - 1} pairs, so the new row has \token{n + 1} numbers.

\lstinputlisting[style=coloredverbatim, caption={Solution for exercise 8, \textit{examples/collections/solutions/pascal.yr}}]{examples/collections/solutions/pascal.yr}
\vspace{-10pt}%
\begin{lstlisting}[style=bashVerb]
[1]
[1, 1]
[1, 2, 1]
[1, 3, 3, 1]
[1, 4, 6, 4, 1]
[1, 5, 10, 10, 5, 1]
[1, 6, 15, 20, 15, 6, 1]
[1, 7, 21, 35, 35, 21, 7, 1]
[1, 8, 28, 56, 70, 56, 28, 8, 1]
[1, 9, 36, 84, 126, 126, 84, 36, 9, 1]
\end{lstlisting}

\subsection*{Exercise 9}

\token{line} returns the letter of three cells when they hold the same one, and
nothing otherwise, an empty cell included. \token{winner} tries each row and
each column in one loop, then the two diagonals.

\lstinputlisting[style=coloredverbatim, caption={Solution for exercise 9, \textit{examples/collections/solutions/tictactoe.yr}}]{examples/collections/solutions/tictactoe.yr}
\vspace{-10pt}%
\begin{lstlisting}[style=bashVerb]
X O X
O X O
. O X
X wins.

O X .
O X X
O . .
O wins.

X O X
X O O
O X X
No winner.
\end{lstlisting}

\subsection*{Exercise 10}

The cell of row \token{i} and column \token{j} is at index \token{i * 3 + j}.
Row \token{i} is therefore made of the indices \token{i * 3} to
\token{i * 3 + 2}, and column \token{j} of the indices \token{j}, \token{3 + j}
and \token{6 + j}: the first column is 0, 3 and 6. The diagonals are 0, 4 and 8,
and 2, 4 and 6. \token{line} is unchanged, and \token{show} starts a new line
after each cell whose index leaves a remainder of 2 when divided by 3, the last
of its row. The program prints the same as that of Exercise~9.

\lstinputlisting[style=coloredverbatim, caption={Solution for exercise 10, \textit{examples/collections/solutions/tictactoe\_flat.yr}}]{examples/collections/solutions/tictactoe_flat.yr}
